\capitulo{3}{Conceptos teóricos}

Algunos conceptos teóricos de \LaTeX{} \footnote{Créditos a los proyectos de Álvaro López Cantero: Configurador de Presupuestos y Roberto Izquierdo Amo: PLQuiz}.

\section{Secciones}

Las secciones se incluyen con el comando section.

\subsection{Subsecciones}

Además de secciones tenemos subsecciones.

\subsubsection{Subsubsecciones}

Y subsecciones. 

\section{Visión por computador}

\subsection{Visión por Computador}
La Visión por Computador es la rama de la Inteligencia Artificial que estudia cómo dotar a un sistema de la capacidad de interpretar imágenes y vídeo, de modo que una matriz de píxeles se transforme en información estructurada y accionable sobre el entorno, como categorías, localizaciones, contornos u otras magnitudes geométricas relevantes. 

Esta disciplina se apoya, por un lado, en la modelización del proceso de formación de imagen y de la geometría de proyección, lo cual permite relacionar el mundo tridimensional con el plano imagen, tal y como se ilustra en la Figura~\ref{fig:Pinhole_camera_model_technical_version}. Por otro lado, operacionaliza ese conocimiento mediante un \textit{pipeline} de procesamiento que convierte progresivamente la señal visual en representaciones más informativas, encadenando transformaciones que filtran, realzan y reorganizan la información hasta hacerla útil para una tarea concreta, como se aprecia en la Figura~\ref{fig:Image_processing_pipeline}. 

En el enfoque contemporáneo, estas transformaciones dejan de diseñarse manualmente y pasan a aprenderse a partir de datos mediante redes profundas, donde capas sucesivas construyen descriptores jerárquicos con creciente carga semántica, un esquema típico se resume en la Figura~\ref{fig:DFVC}. Sobre estas representaciones, el sistema ejecuta tareas como clasificación, detección o segmentación, y en particular la segmentación semántica permite asignar una etiqueta a cada región de la imagen preservando estructura espacial, tal y como ejemplifica el proceso de la Figura~\ref{fig:Image-segmentation-example}
y la Figura~\ref{fig:Image-segmentation-example-segmented}, conectando de forma natural con los fundamentos que se desarrollan en las secciones posteriores.

\subsection{Conceptos teóricos de la Segmentación Semántica}
En esta sección se introducen los fundamentos teóricos que sustentan la segmentación semántica moderna en visión por computador, abordándola desde la perspectiva del aprendizaje profundo y de las arquitecturas que permiten realizar predicción densa a nivel de píxel. Se describen los principios que rigen el uso de redes neuronales profundas, haciendo especial énfasis en el papel de las estructuras ultra-profundas y de los mecanismos que facilitan su entrenamiento, como las conexiones residuales y las \textit{skip connections}. Asimismo, se analizan los componentes operativos esenciales de los modelos de codificador-decodificador, incluyendo conceptos clave como \textit{kernels}, \textit{stride}, \textit{downsampling} y \textit{upsampling}, los cuales determinan el compromiso entre resolución espacial y riqueza semántica. Finalmente, se presentan las dos familias arquitectónicas predominantes en este ámbito, las redes neuronales convolucionales y los \textit{Vision Transformers}, destacando sus principios de funcionamiento y su relevancia en tareas de segmentación semántica de alta precisión.


\subsubsection{Aprendizaje Profundo}
El Aprendizaje Profundo es un paradigma de machine learning basado en el uso de redes neuronales convolucionales (CNN) o arquitecturas alternativas como los \textit{Transformers}, que han surgido como el enfoque dominante para el reconocimiento visual~\cite{he2016resnet, huang2018densenet}. 
Funciona de una forma que, una clase de hipótesis de altísima capacidad, realizada por composiciones profundas de operadores paramétricos y no linealidades puntuales, aprende representaciones jerárquicas de los datos optimizando un funcional de riesgo empírico.

El éxito del Aprendizaje Profundo se atribuye a su capacidad para trabajar con arquitecturas muy grandes y profundas, con millones de parámetros, logrando avances significativos en el reconocimiento de imágenes~\cite{ronneberger2015unet, diezpastora2025trails, neloy2024autoencoders, he2016resnet}. 
Este resuelve la necesidad de los sistemas de visión de extraer la representación de características más efectiva a partir de datos complejos~\cite{neloy2024autoencoders}, automatizando lo que antes requería el diseño manual de características.

\subsubsection{Estructuras ultra-profundas}
En el marco del Aprendizaje Profundo, la profundidad de las representaciones y de la red es de importancia crucial para muchas tareas de reconocimiento visual~\cite{he2016resnet, huang2018densenet}. Las arquitecturas modernas han superado significativamente las redes iniciales, con modelos que han rebasado la barrera de las cien capas~\cite{huang2018densenet}. Estas estructuras ultra-profundas buscan enriquecer los niveles de características mediante el aumento del número de capas apiladas~\cite{he2016resnet, huang2018densenet}.

Sin embargo, el aumento de la profundidad introduce un desafío conocido como el problema de la degradación. Esta ocurre cuando, al aumentar las capas apiladas, la precisión del entrenamiento se deteriora. Teóricamente, un modelo más profundo no debería tener un error de entrenamiento mayor que su contraparte menos profunda. El hecho de que esto suceda indica que los algoritmos de optimización actuales tienen dificultades para encontrar soluciones que optimicen estas estructuras ultra-profundas~\cite{he2016resnet}. 

Esto debido a que, a medida que la información o el gradiente pasan por muchas capas, pueden desvanecerse antes de alcanzar el inicio o el final de la red~\cite{huang2018densenet}.

\subsubsection{Predición Densa (Dense prediction)}
La predicción densa es una formulación de inferencia estructurada en la que el sistema estima, para cada ubicación del dominio espacial de la entrada, una variable de salida perteneciente a un espacio de etiquetas o a un espacio continuo~\cite{zhang2022constructionautomation}. En visión por computador abarca tanto clasificación por elemento, como regresión por elemento, dando lugar a campos de salida con fuerte dependencia local y coherencia geométrica~\cite{zhao2017pspnet, lin2017fpn}. 

Desde el punto de vista inductivo, estas tareas se benefician de la \textit{equivarianza traslacional} propia de arquitecturas totalmente convolucionales, del acoplamiento multi-escala para capturar contexto y del control explícito del \textit{output stride} para balancear resolución y campo receptivo~\cite{ liu2022convnext, khan2022transformersvision}. La naturaleza densa impone consideraciones de regularidad de frontera, conservación de detalles de alta frecuencia, manejo del desbalanceo de clases y robustez frente a los problemas introducidos por downsampling~\cite{cai2023efficientvit, xiao2018upernet, ji2022automation}.

Operativamente, la predicción densa se implementa con \textit{encoders–decoders} convolucionales que combinan \textit{downsampling} para agregar contexto y \textit{upsampling} para reconstruir resolución, asistidos por \textit{skip connections} por concatenación para preservar información de detalle~\cite{ronneberger2015unet, bala2025visiontranscoder, badrinarayanan2016segnet}.

\subsubsection{Segmentación Semántica}
La segmentación semántica es una tarea de predicción densa en Visión por Computadora cuyo objetivo es la clasificación a nivel de píxel~\cite{bala2025visiontranscoder, ji2022automation}, mediante la asignación de una etiqueta de clase semántica a cada píxel de una imagen de entrada. A diferencia de tareas como la clasificación, la segmentación semántica proporciona una comprensión completa del escenario, prediciendo así la categoría, ubicación y forma de cada elemento.

Este proceso típicamente se implementa a través de una arquitectura de red neuronal profunda de codificador-decodificador, llamada \textit{encoder-decoder} de manera tradicional~\cite{ji2022automation, neloy2024autoencoders, bala2025visiontranscoder}.

\begin{enumerate}
	\item \textbf{Codificación mediante encoder:} 
	El codificador, a menudo basado en una Red Neuronal Convolucional (CNN)~\cite{visionencoder2025} o un \textit{Vision Transformer (ViT)}~\cite{bala2025visiontranscoder}, toma la imagen de entrada y transforma los datos de esta en una representación de características de alto nivel mediante una reducción de la resolución espacial a través de diversas operaciones como podrían ser el \textit{max-pooling} o el uso de capas convolucionales, obteniendo así como resultado unas capas de características de baja resolución pero semánticamente ricas.
	
	\item \textbf{Decodificación por decoder:} El decodificador opera con las características comprimidas del codificador para reconstruir la predicción final. Esto lo hará mediante la utilización de operaciones de sobremuestreo, a menudo implementadas mediante \textit{upsampling}~\cite{badrinarayanan2016segnet, lin2017fpn} o mediante el uso de índices de \textit{max-pooling}~\cite{neloy2024autoencoders, zhang2022constructionautomation}, agrandando la entrada para extraer características de alta resolución~\cite{neloy2024autoencoders, ji2022automation, badrinarayanan2016segnet}.
	
	\item \textbf{Salida y Optimización:} La salida final es un mapa de segmentación que se alimenta a una capa \textit{Softmax} para asignar probabilidades de clasificación a cada píxel de forma independiente~\cite{ji2022automation, bala2025visiontranscoder, neloy2024autoencoders}, generando una predicción densa~\cite{neloy2024autoencoders, bala2025visiontranscoder}.
\end{enumerate}

\subsubsection{Kernel}
En el ámbito de las redes neuronales, un \textit{kernel} es un conjunto de filtros con pesos aprendibles que constituyen el componente fundamental de las capas convolucionales para convertir datos de entrada en mapas de características~\cite{ji2022automation}. Operativamente, el \textit{kernel} se desliza sobre la entrada utilizando un \textit{stride}, realizando una operación matemática que produce una salida espacial encargada de capturar jerarquías de información, desde bordes simples hasta formas complejas~\cite{he2016resnet, ji2022automation, liu2022convnext}. Esta estrategia es intrínseca al procesamiento visual, permitiendo que las computaciones se compartan y que el sistema sea eficiente al manejar imágenes de alta resolución~\cite{liu2022convnext}.

El diseño del \textit{kernel}, es una decisión arquitectónica crítica que determina el campo receptivo; por ejemplo, mientras que los \textit{kernels} de $3 \times 3$ son el estándar por su eficiencia en \textit{hardware}, modelos modernos han demostrado que los \textit{kernels} de $7 \times 7$ pueden capturar un contexto global mucho más amplio~\cite{liu2022convnext}. 

Además de la extracción de características, el concepto de \textit{kernel} se extiende a funciones matemáticas de comparación, como el \textit{Central Kernel Alignment (CKA)}, que utiliza un \textit{kernel} lineal para medir la similitud de representación entre los espacios de embebido latente de diferentes arquitecturas~\cite{iana2024peeling}.

\subsubsection{Stride}
El \textit{stride} es un parámetro crítico que define la magnitud del desplazamiento de un \textit{kernel} a medida que recorre los datos de entrada en una red neuronal~\cite{ji2022automation, dumoulin2018guide}. En las capas del codificador, el uso de un \textit{stride} mayor a uno constituye una técnica de \textit{downsampling} que reduce progresivamente la resolución espacial de los mapas de características~\cite{badrinarayanan2016segnet, ronneberger2015unet}. Esta operación permite que los píxeles resultantes abarquen un contexto de imagen de entrada más amplio, facilitando la extracción de características semánticas de alto nivel y mejorando la invarianza traslacional del modelo~\cite{ji2022automation, badrinarayanan2016segnet, liu2022convnext}.

En el decodificador, el \textit{stride} es esencial para las operaciones de \textit{upsampling}, donde se utiliza para expandir representaciones comprimidas hacia mapas de alta resolución~\cite{ji2022automation, badrinarayanan2016segnet}. El tamaño final de la salida depende de la interacción matemática entre las dimensiones de entrada, el \textit{padding}, el tamaño del \textit{kernel} y el valor del \textit{stride} aplicado~\cite{ji2022automation, tan2019efficientnet}. La gestión precisa de este paso es vital para reconstruir con fidelidad la información espacial y asegurar una delineación precisa de los límites de los objetos en tareas de segmentación densa~\cite{ronneberger2015unet, howard2019mobilenetv3, bala2025visiontranscoder}.

\subsubsection{Transformer}
La arquitectura del \textit{Transformer} es un modelo de procesamiento de secuencias basado íntegramente en mecanismos de autoatención, eliminando la necesidad de recurrencia o convoluciones para capturar dependencias~\cite{vaswani2017attention, khan2022transformersvision}. A diferencia de las redes recurrentes que procesan elementos de forma secuencial, el \textit{Transformer} permite el procesamiento en paralelo de toda la secuencia, lo que facilita el modelado de relaciones de largo alcance y mejora significativamente la escalabilidad en grandes conjuntos de datos~\cite{khan2022transformersvision, mamba2023linear}. 

Su estructura estándar se compone de un codificador y un decodificador que apilan múltiples bloques idénticos, cada uno equipado con capas de atención de múltiples cabezales y redes de alimentación \textit{feed-forward}~\cite{vaswani2017attention, bala2025visiontranscoder}.

\subsubsection{Mecanismos Residuales}
Los mecanismos residuales son esquemas de conexión aditiva que introducen \textit{skip connections} a través de los bloques de una red profunda, de modo que la transformación aprendida se formula como una corrección respecto de la señal que fluye sin distorsión por el atajo~\cite{he2016resnet}. 

Esta arquitectura estabiliza la propagación hacia delante y hacia atrás, mejora el acondicionamiento efectivo del problema de optimización y habilita redes sustancialmente más profundas; cuando existe desajuste de dimensionalidad, la ruta de identidad se implementa mediante proyecciones, mientras que otras variantes con saltos de identidad estrictos refuerzan el transporte directo de información y gradientes.

En conjunto, los mecanismos residuales facilitan el aprendizaje de funciones de alta complejidad al mantener rutas de baja impedancia para el flujo de señales~\cite{srivastava2015highway}.

\subsubsection{Skip Connections}
Conexiones encargadas de mitigar la pérdida de información espacial que ocurre durante el \textit{downsampling} en el codificador~\cite{bala2025visiontranscoder}. Transfieren características de alta resolución desde las capas tempranas del codificador a las capas correspondientes del decodificador, fusionando las características codificadas con los detalles espaciales finos, siendo estos los límites y contornos, del inicio de la red~\cite{badrinarayanan2016segnet, ronneberger2015unet}. 

Para lograr esto, se emplea una capa de concatenación que fusiona la salida de las capas convolucionales del \textit{encoder} con las capas de deconvolución del \textit{decoder}~\cite{ronneberger2015unet, ji2022automation}.

\subsubsection{Upsampling}
El \textit{upsampling}, también denominado deconvolución, es la función principal del decodificador, definida como la operación de mapeo que transforma representaciones de características comprimidas de baja resolución en mapas de características de alta resolución espacial para la clasificación a nivel de píxel~\cite{badrinarayanan2016segnet, ji2022automation}. 

Esto puede implementarse mediante interpolación, convolución transpuesta o decodificadores MLP~\cite{ji2022automation, cai2023efficientvit}. La implementación también puede optar por utilizar métodos sin aprendizaje, aunque esto puede resultar en un rendimiento inferior en comparación con los métodos que sí aprenden del entorno~\cite{badrinarayanan2016segnet}.

Para preservar los contornos finos y garantizar la localización precisa, el decodificador se basa en la fusión de características de niveles previos del codificador mediante \textit{skip connections}~\cite{ronneberger2015unet, badrinarayanan2016segnet, bala2025visiontranscoder}.

\subsubsection{Downsampling}
El \textit{downsampling} es una técnica fundamental utilizada en las redes neuronales convolucionales y en el \textit{encoder} de las arquitecturas de codificador-decodificador para transformar la información de entrada en una representación de características reducida de alto nivel~\cite{ji2022automation, badrinarayanan2016segnet}. 

El propósito central de esta operación es doble: por un lado, aliviar la carga computacional y prevenir el \textit{overfitting}; por otro lado, lograr una mayor \textit{invarianza traslacional} sobre pequeños desplazamientos espaciales en la imagen de entrada, lo cual robustece la clasificación. Esto se realiza a costa de una pérdida progresiva de la resolución espacial en los mapas de características, lo cual reduce el detalle de los límites, un factor que debe compensarse en la fase de decodificación~\cite{badrinarayanan2016segnet, ronneberger2015unet, bala2025visiontranscoder}.

El mecanismo principal para lograr el submuestreo es la operación de \textit{max-pooling}, que es el núcleo del \textit{downsampling} y se utiliza para disminuir el tamaño del mapa de características~\cite{ji2022automation, ronneberger2015unet}. Lo que hace para llevarlo a cabo es dividir dicho mapa de entrada en pequeñas regiones o ventanas. Tras ello, extrae el valor máximo dentro de una de las ventanas, utilizándolo como el único valor representativo para esa área en el nuevo mapa de características reducido. Finalmente, la ventana se desliza a lo largo de la entrada con un \textit{stride} predefinido, permitiendo que cada píxel resultante abarque un contexto de imagen de entrada más amplio a medida que se reduce la resolución~\cite{zhang2022constructionautomation, ji2022automation, badrinarayanan2016segnet}. 

Aunque las sucesivas capas de \textit{max-pooling} y submuestreo pueden lograr una mayor invarianza para una clasificación robusta, la contraparte es una pérdida progresiva de la resolución espacial en los mapas de características. Esta pérdida resulta en una representación cada vez más inexacta en cuanto al detalle de los límites de los objetos, lo cual es perjudicial para tareas de predicción densa, como la segmentación, donde la delineación precisa es vital~\cite{badrinarayanan2016segnet, bala2025visiontranscoder}. Sin embargo, el \textit{downsampling} tiene un beneficio inherente: logra que cada píxel en el mapa de características posterior abarque un contexto de imagen de entrada más grande.


\subsubsection{Red Neuronal Convolucional (CNN)}
Una Red Neuronal Convolucional es una arquitectura de aprendizaje profundo que se ha convertido en el enfoque dominante para el reconocimiento visual de objetos~\cite{huang2018densenet, he2016resnet}. Implementa sesgos inductivos inherentes, como la covarianza traslacional, para la extracción eficiente de características jerárquicas, integrando naturalmente características de nivel bajo, medio y alto~\cite{liu2022convnext}. En un contexto de predicción densa, como la segmentación semántica, la CNN se extiende típicamente en un marco de codificador-decodificador para realizar una clasificación a nivel de píxel~\cite{zhang2022constructionautomation, ji2022automation, badrinarayanan2016segnet}.

El funcionamiento en la segmentación se basa en este marco bifásico de codificación y decodificación:

\begin{enumerate}
\item \textbf{Capas Convolucionales, conocidas como: Convolutional Layers}

Las capas convolucionales son el núcleo de la CNN y componen la ruta principal del codificador~\cite{ji2022automation, badrinarayanan2016segnet}. En estas, se realiza la operación de convolución, que utiliza un conjunto de \textit{kernels} con pesos entrenables para convertir los datos de entrada en mapas de características~\cite{ji2022automation}. 

El \textit{kernel} se desliza sobre la entrada con un \textit{stride} para producir la salida espacial~\cite{he2016resnet}. Mediante este proceso de \textit{downsampling} codifica la información de entrada en una representación de características de alto nivel downsampleada~\cite{ji2022automation, badrinarayanan2016segnet}.

\item \textbf{Capas Deconvolucionales, también conocidas como: Deconvolutional Layers}

Las capas deconvolucionales son el mecanismo de \textit{upsampling} más común utilizado en la ruta del decodificador~\cite{zhang2022constructionautomation, badrinarayanan2016segnet}. Mediante el uso de esta técnica, también referida como convolución transpuesta, consiguen mapear los mapas de características comprimidos de baja resolución a mapas de características densos que se aproximan a la resolución de entrada original para la clasificación por elemento~\cite{badrinarayanan2016segnet, ji2022automation}. 

El desafío al que se deben enfrentar estas capas es que la red puede tender a ignorar los detalles debido a la naturaleza de la amplificación, lo que conlleva una salida insatisfactoria~\cite{ji2022automation}.

\item \textbf{Capa de Concatenación, comunmente: Concatenation Layer}

La capa de concatenación es un componente de fusión de características utilizado en el decodificador para mejorar la precisión espacial~\cite{ ronneberger2015unet, badrinarayanan2016segnet, bala2025visiontranscoder}. Para lograr este objetivo, fusiona las características de salida de las capas convolucionales del codificador con las características de las capas deconvolucionales del decodificador.

Dicha fusión se logra a través del uso de \textit{skip connections}, que copian los mapas de características de alta resolución y bajo nivel semántico desde el codificador hacia el decodificador~\cite{ ronneberger2015unet, badrinarayanan2016segnet, ji2022automation}. La concatenación expande estos mapas y, al combinar las características espaciales detalladas con las características semánticas restauradas, asegura una delineación precisa de los límites en la predicción final~\cite{ ronneberger2015unet, badrinarayanan2016segnet, bala2025visiontranscoder}.

\end{enumerate}

\subsubsection{Vision Transformer (ViT)}
La arquitectura del \textit{Vision Transformer (ViT)} se erige como el paradigma dominante después de las CNN en tareas de reconocimiento visual~\cite{liu2022convnext, khan2022transformersvision, bala2025visiontranscoder}. Opera al inicio con una fase de tokenización no superpuesta, donde la imagen de entrada es segmentada en parches, aplanados, y proyectados linealmente mediante una transformación a un espacio embebido~\cite{bala2025visiontranscoder}. A estos \textit{embeddings} se les adjuntan codificaciones posicionales para conservar la información espacial~\cite{bala2025visiontranscoder}.

Esta secuencia de \textit{tokens} enriquecida se introduce en un \textit{encoder} de \textit{Transformer} que consiste en múltiples bloques idénticos, equipados con el mecanismo de \textit{Multi-Head Self-Attention}. El \textit{self-attention} permite que cada \textit{token} capture relaciones de largo alcance y el contexto global con todos los demás \textit{tokens}, facilitando una interpretación más inclusiva de patrones intrincados en los datos~\cite{khan2022transformersvision, bala2025visiontranscoder}.

En el contexto de la segmentación densa, el \textit{encoder ViT} finaliza su proceso entregando características codificadas de bajo nivel espacial pero semánticamente ricas. La arquitectura completa luego emplea un decodificador para reconstruir el mapa de segmentación~\cite{bala2025visiontranscoder}. Este decodificador típicamente incluye capas de \textit{upsampling} para aumentar la resolución espacial de las características comprimidas y utiliza \textit{skip connections} para transferir características de alta resolución desde el \textit{encoder} y así reconstruir el mapa de segmentación con alta fidelidad~\cite{bala2025visiontranscoder}. 

La salida final se genera a través de una capa convolucional seguida de una activación \textit{Softmax}, proporcionando un mapa de probabilidades detallado para la clasificación a nivel de píxel~\cite{zhang2022constructionautomation, bala2025visiontranscoder}. 

\subsection{Arquitecturas de Segmentación Semántica}
En esta sección se presentan las principales arquitecturas de segmentación semántica empleadas en el ámbito de la visión por computador, las cuales materializan los principios teóricos descritos anteriormente mediante diseños de red específicos orientados a la predicción densa. Se analizan arquitecturas representativas basadas en esquemas de codificador-decodificador y en la agregación jerárquica de características, como U-Net, Feature Pyramid Network y UPerNet, así como enfoques más recientes que incorporan mecanismos de atención global mediante \textit{Transformers}, como SegFormer. Asimismo, se incluye PSPNet, una arquitectura pionera en la integración explícita de contexto multi-escala a través de pirámides de agrupamiento. Cada una de estas arquitecturas introduce estrategias diferenciadas para capturar información semántica y preservar la coherencia espacial, constituyendo referentes fundamentales en el desarrollo y la evolución de los modelos de segmentación semántica.


\subsubsection{U-Net}
La red neuronal convolucional U-Net constituye un paradigma fundamental en el ámbito de la segmentación densa debido a su capacidad para operar eficientemente con conjuntos de datos de entrenamiento limitados~\cite{ronneberger2015unet}. Esta arquitectura presenta una topología simétrica en forma de U que integra un codificador y un decodificador~\cite{ronneberger2015unet, Dimitrovski2024}. El codificador aplica sucesivamente convoluciones seguidas de unidades de activación lineal rectificada y operaciones de \textit{max-pooling} para la extracción de características contextuales y la reducción de la dimensionalidad espacial~\cite{ronneberger2015unet, ji2022automation}. Por su parte, el decodificador realiza procesos de \textit{upsampling} con el fin de recuperar la resolución original de la imagen~\cite{ronneberger2015unet, badrinarayanan2016segnet}. 

El componente diferencial de este modelo reside en la implementación de conexiones de salto que concatenan mapas de características de alta resolución del codificador con las capas correspondientes del decodificador~\cite{ronneberger2015unet, diezpastora2025trails}. Dicho mecanismo permite la preservación de detalles espaciales que se degradan habitualmente durante el submuestreo, garantizando una localización precisa de los límites de los objetos en tareas críticas como la imagen médica o la teledetección~\cite{Dimitrovski2024}.

\subsubsection{Feature Pyramid Network}
La arquitectura Feature Pyramid Network constituye un extractor de características de propósito general diseñado para explotar la jerarquía piramidal inherente a las redes neuronales convolucionales profundas~\cite{lin2017fpn}. Este framework implementa una ruta descendente con conexiones laterales que permite la propagación de información semántica de alto nivel hacia mapas de características de mayor resolución espacial~\cite{lin2017fpn, ji2022automation}. 

Operacionalmente, el sistema fusiona representaciones semánticamente robustas provenientes de capas profundas con mapas de características semánticamente débiles pero espacialmente precisos de etapas tempranas~\cite{lin2017fpn, xiao2018upernet}. La importancia de este diseño radica en su capacidad para generar una representación multiescala donde todos los niveles poseen una fuerte carga semántica, facilitando la detección y segmentación de objetos en un rango extenso de escalas sin incurrir en costes computacionales significativos~\cite{lin2017fpn}.


\subsubsection{UPerNet}
El modelo \textit{Unified Perceptual Parsing Network}, denominado \textit{UPerNet}, se erige sobre la arquitectura de redes de pirámide de características para resolver la segmentación de múltiples conceptos visuales de forma simultánea~\cite{xiao2018upernet}. Este diseño estructural adopta el esquema de una \textit{Feature Pyramid Network} y lo potencia mediante la integración de un \textit{Pyramid Pooling Module} sobre la capa de características más profunda del codificador~\cite{xiao2018upernet, zhao2017pspnet}. El funcionamiento del sistema se basa en la fusión jerárquica de representaciones provenientes de diversas etapas de la red, donde los mapas de niveles superiores aportan información semántica global y los niveles inferiores proporcionan detalles geométricos finos~\cite{xiao2018upernet, lin2017fpn}. 

Mediante la interpolación bilineal y la concatenación de estos mapas de características multiescala, UPerNet logra una proyección no lineal eficiente hacia un espacio de alta dimensionalidad que captura relaciones espaciales complejas~\cite{xiao2018upernet}. Su importancia en el estado del arte radica en la capacidad para generar mapas de segmentación semánticamente ricos sin incurrir en la elevada carga computacional de las convoluciones dilatadas tradicionales, permitiendo la identificación de objetos, sus partes y materiales bajo un único proceso de inferencia unificado~\cite{xiao2018upernet, lin2017fpn}.

\subsubsection{SegFormer}
SegFormer representa un paradigma avanzado de segmentación semántica que unifica codificadores basados en transformadores jerárquicos con decodificadores ligeros de perceptrón multicapa~\cite{xie2021segformer, Dimitrovski2024}. A diferencia de los Vision Transformers convencionales, esta arquitectura introduce un codificador libre de codificaciones posicionales mediante el uso de convoluciones superpuestas, lo que permite inferir sobre resoluciones de imagen variables sin degradación del rendimiento~\cite{xie2021segformer, cai2023efficientvit}. 

El funcionamiento del sistema se fundamenta en la generación de características multiescala que capturan simultáneamente dependencias de largo alcance y detalles locales, enviando estas representaciones a un decodificador All-MLP que agrega eficientemente la información contextual~\cite{xie2021segformer, Dimitrovski2024}. La relevancia de SegFormer en el estado del arte se debe a su alta eficiencia computacional y robustez, logrando resultados competitivos en tareas de predicción densa de alta resolución al mitigar la complejidad cuadrática inherente a los mecanismos de autoatención estándar~\cite{xie2021segformer, cai2023efficientvit}.

\subsubsection{PSPNet (Pyramid Scene Parsing Network)}
El framework \textit{Pyramid Scene Parsing Network} constituye un avance fundamental en la segmentación de escenas complejas al introducir mecanismos para la explotación exhaustiva de la información de contexto global~\cite{zhao2017pspnet}. 

Esta arquitectura surge como respuesta a la limitación intrínseca de las redes completamente convolucionales tradicionales, donde el campo receptivo empírico resulta insuficiente para capturar dependencias espaciales de largo alcance en capas profundas~\cite{zhao2017pspnet, xiao2018upernet}. La pieza central del modelo es el \textit{Pyramid Pooling Module} (PPM), el cual recolecta representaciones de subregiones en múltiples escalas para construir una jerarquía de prioridades globales que guían la clasificación a nivel de píxel~\cite{zhao2017pspnet}.

Operacionalmente, el PPM divide el mapa de características generado por un codificador basado en bloques residuales en cuatro niveles de granularidad distintos, aplicando operaciones de \textit{pooling} con tamaños de bin de $1\times 1$, $2\times 2$, $3\times 3$ y $6\times 6$~\cite{zhao2017pspnet, he2016resnet}. A continuación, cada nivel se somete a una convolución de $1\times 1$ para realizar una reducción de dimensionalidad y se remuestrea mediante interpolación bilineal hasta alcanzar la resolución original de la etapa~\cite{zhao2017pspnet, lin2017fpn}. Estos mapas enriquecidos se concatenan con el flujo de información original para formar una representación final que encapsula tanto el detalle geométrico local como la coherencia semántica global, mitigando errores comunes en la identificación de objetos con apariencias similares pero contextos dispares~\cite{zhao2017pspnet, diezpastora2025trails}.

Para garantizar la viabilidad de la optimización en estructuras ultra-profundas, el entrenamiento de PSPNet se apoya en una estrategia de aprendizaje profundamente supervisado mediante el uso de una pérdida auxiliar~\cite{zhao2017pspnet}. Este mecanismo aplica una función de coste adicional en etapas intermedias del codificador, típicamente tras el cuarto bloque residual, lo que facilita la propagación del gradiente y actúa como un regularizador durante la fase de convergencia~\cite{zhao2017pspnet, he2016resnet}.

\section{Referencias}

Las referencias se incluyen en el texto usando cite~\cite{wiki:latex}. Para citar webs, artículos o libros~\cite{koza92}, si se desean citar más de uno en el mismo lugar~\cite{bortolot2005, koza92}.


\section{Imágenes}

Se pueden incluir imágenes con los comandos standard de \LaTeX, pero esta plantilla dispone de comandos propios como por ejemplo el siguiente:

\imagen{escudoInfor}{Autómata para una expresión vacía}{.5}

\imagen{Pinhole_camera_model_technical_version}{Modelo de cámara \textit{pinhole} como abstracción geométrica del proceso de formación de imagen. Fuente: Olaf Peters, licencia CC BY-SA 4.0.}{1}

\imagen{Image_processing_pipeline}{Ejemplo de \textit{pipeline} de procesamiento de imagen, donde la entrada se transforma mediante etapas sucesivas. Fuente: Nzebro12, licencia CC BY-SA 4.0.}{1}

\imagen{DFVC}{Diagrama que detalla el flujo correspondiente a un sistema que emplea el uso de visión por computador}{1}

\imagen{Image-segmentation-example}{Ejemplo de segmentación, donde distintas regiones quedan diferenciadas para facilitar interpretación a nivel de escena (imagen original). Fuente: MartinThoma, licencia CC0 1.0.}{1}

\imagen{Image-segmentation-example-segmented}{Ejemplo de segmentación, donde distintas regiones quedan diferenciadas para facilitar interpretación a nivel de escena (imagen segmentada). Fuente: MartinThoma, licencia CC0 1.0.}{1}


\section{Listas de items}

Existen tres posibilidades:

\begin{itemize}
	\item primer item.
	\item segundo item.
\end{itemize}

\begin{enumerate}
	\item primer item.
	\item segundo item.
\end{enumerate}

\begin{description}
	\item[Primer item] más información sobre el primer item.
	\item[Segundo item] más información sobre el segundo item.
\end{description}
	
\begin{itemize}
\item 
\end{itemize}

\section{Tablas}

Igualmente se pueden usar los comandos específicos de \LaTeX o bien usar alguno de los comandos de la plantilla.

\tablaSmall{Herramientas y tecnologías utilizadas en cada parte del proyecto}{l c c c c}{herramientasportipodeuso}
{ \multicolumn{1}{l}{Herramientas} & App AngularJS & API REST & BD & Memoria \\}{ 
HTML5 & X & & &\\
CSS3 & X & & &\\
BOOTSTRAP & X & & &\\
JavaScript & X & & &\\
AngularJS & X & & &\\
Bower & X & & &\\
PHP & & X & &\\
Karma + Jasmine & X & & &\\
Slim framework & & X & &\\
Idiorm & & X & &\\
Composer & & X & &\\
JSON & X & X & &\\
PhpStorm & X & X & &\\
MySQL & & & X &\\
PhpMyAdmin & & & X &\\
Git + BitBucket & X & X & X & X\\
Mik\TeX{} & & & & X\\
\TeX{}Maker & & & & X\\
Astah & & & & X\\
Balsamiq Mockups & X & & &\\
VersionOne & X & X & X & X\\
} 
